\documentclass[12pt,a4paper]{article}
\usepackage[utf8]{inputenc}
\usepackage{amsmath,amssymb,amsthm}
\usepackage{physics}
\usepackage{booktabs}
\usepackage{graphicx}
\usepackage{hyperref}
\usepackage[margin=2.5cm]{geometry}
\usepackage{enumitem}
\usepackage{xcolor}
\usepackage{tcolorbox}

\newtheorem{theorem}{Theorem}
\newtheorem{proposition}[theorem]{Proposition}
\newtheorem{corollary}[theorem]{Corollary}
\theoremstyle{definition}
\newtheorem{definition}[theorem]{Definition}
\theoremstyle{remark}
\newtheorem{remark}[theorem]{Remark}

\tcbset{colback=blue!5!white, colframe=blue!75!black, fonttitle=\bfseries}

\title{The Cooper Pair Mass Anomaly as a Propulsion Mechanism:\\
\large DFD Prediction $\delta = \sqrt{3}\,\alpha^2$, Experimental Constraints,\\
and a Tabletop Test of Gravitational Coupling}
\author{G.T.\ Alcock\\
\textit{Density Field Dynamics Research Programme}}
\date{March 2026}

\begin{document}
\maketitle

\begin{abstract}
We analyse the 36-year-old anomaly in the Cooper pair effective mass measured by
Tate, Cabrera, Felch, and Anderson (1989, 1990) in rotating niobium superconductors.
The measured ratio $m^*/2m_e = 1.000\,084(21)$ disagrees with BCS theory by
$92 \pm 21$~ppm.  Density Field Dynamics (DFD) predicts \emph{exactly} this anomaly:
$\delta = \sqrt{3}\,\alpha^2 = 92.23$~ppm, arising from the quintet exchange channel
$|5\rangle \in S^2(V^*)$ of the $A_5$ microsector.  The critical question for propulsion
is whether this anomalous mass \emph{gravitates}.  We show that the Ross~et~al.\ (2024)
E\"ot-Wash cryogenic torsion balance constrains the gravitational coupling parameter
$\lambda_{\rm quintet}$ only to $|1 - \lambda_{\rm quintet}| \leq 10$, leaving the
propulsion-relevant regime $\lambda_{\rm quintet} = 0$ completely unconstrained.  For a
10~kg niobium sample cooled below $T_c = 9.25$~K, the predicted weight anomaly is
$\Delta W = 3.44$~mN (350~mg$_f$) if $\lambda_{\rm quintet} = 0$---detectable at
signal-to-noise $> 35{,}000\!:\!1$ on a standard 0.01~mg analytical balance.  We design
a complete experimental protocol using thermal cycling through $T_c$, estimate all
systematic errors, and show this test can be performed in any university cryogenics
laboratory for less than \$50{,}000.  This mechanism is fundamentally different from all
previous propulsion proposals: it does not attempt to \emph{create} $\psi$-field gradients
but instead modifies how an existing gravitational field couples to a quantum condensate.
\end{abstract}

\tableofcontents

%=============================================================================
\section{Introduction}\label{sec:intro}
%=============================================================================

In 1989, Tate, Cabrera, Felch, and Anderson published a precision measurement of
the Cooper pair effective mass in niobium using the London moment of a rotating
superconductor~\cite{Tate1989,Tate1990}.  They found:
\begin{equation}\label{eq:tate}
\frac{m^*}{2m_e} = 1.000\,084 \pm 0.000\,021,
\end{equation}
a result that disagrees with the BCS theoretical prediction of $m^*/2m_e = 0.999\,992$
by approximately 92~ppm at $4.4\sigma$ significance.

This anomaly has remained unexplained for 36 years.  Various theoretical proposals
have been advanced---gravitomagnetic fields in rotating superconductors~\cite{Tajmar2006},
modifications to the canonical momentum including gravitational
terms~\cite{DeWitt1966}---but none has provided a satisfactory quantitative explanation
from first principles.

Density Field Dynamics (DFD) provides exactly such an explanation.  The $A_5$
microsector, which governs the internal structure of fermionic states through the
channel Hilbert space $\mathcal{H}_{\rm ch} = \mathbb{C}[A_5]$ with
$\dim = |A_5| = 60$, contains a quintet representation $|5\rangle$ in $S^2(V^*)$
that is accessible only to \emph{paired} electrons.  The exchange energy in this
channel contributes:
\begin{equation}\label{eq:delta_DFD}
\boxed{\delta = \sqrt{3}\,\alpha^2 = 9.223 \times 10^{-5} = 92.23~\text{ppm}}
\end{equation}
to the effective Cooper pair mass.  This is within $0.4\sigma$ of the Tate
measurement.

The propulsion question reduces to a single unknown: the gravitational coupling
parameter $\lambda_{\rm quintet} \in [0,1]$, which determines whether the quintet
exchange energy sources the DFD scalar field $\psi$ (i.e., whether it gravitates).

%=============================================================================
\section{The Tate Experiment and Subsequent Measurements}\label{sec:tate}
%=============================================================================

\subsection{Experimental Method}

The Tate experiment measures the London moment---the magnetic field generated by
a rotating superconductor.  When a superconductor rotates with angular velocity
$\Omega$, the Cooper pairs lag behind the ionic lattice, producing a magnetic field:
\begin{equation}\label{eq:london}
\mathbf{B}_L = -\frac{2m^*}{e}\,\boldsymbol{\Omega},
\end{equation}
where $m^*$ is the effective mass per electron in the Cooper pair and $e$ is the
electron charge.  By measuring $B_L$ with a SQUID magnetometer and independently
measuring $\Omega$, one obtains $m^*$.

More precisely, Tate et al.\ measured the ratio $h/m'$ where $m' = 2m^*$ is
the observable Cooper pair mass, using a rotating superconducting niobium ring.
The magnetic flux through the ring vanishes for a discrete set of rotation
frequencies equally spaced by $\Delta\nu$, giving:
\begin{equation}
\frac{h}{m'} = 2S\,\Delta\nu,
\end{equation}
where $S$ is the area bounded by the ring.  This allowed a determination of
$m'/2m_e$ with 5~ppm precision and 21~ppm accuracy.

\subsection{Measurement Results}

\begin{table}[h]
\centering
\renewcommand{\arraystretch}{1.4}
\begin{tabular}{lccc}
\toprule
\textbf{Group} & \textbf{Year} & \textbf{$m^*/2m_e$} & \textbf{Anomaly (ppm)} \\
\midrule
Tate et al.\ (PRL) & 1989 & $1.000\,084(21)$ & $84 \pm 21$ \\
Tate et al.\ (PRB) & 1990 & $1.000\,084(21)$ & $84 \pm 21$ \\
Hoang et al. & 2019 & (multiple SC) & $18 \pm 2.1$ (Nb) \\
\midrule
BCS theory & --- & $0.999\,992$ & $-8$ \\
\textbf{DFD prediction} & --- & $1.000\,092$ & $\mathbf{+92.23}$ \\
\bottomrule
\end{tabular}
\caption{Measurements and predictions for the Cooper pair effective mass in niobium.
The anomaly column gives $m^*/2m_e - 1$ in parts per million.  The DFD prediction
$\delta = \sqrt{3}\,\alpha^2$ is computed from the $A_5$ microsector quintet channel.
Note: the Tate anomaly \emph{above BCS prediction} is $84 - (-8) = 92$~ppm.}
\label{tab:measurements}
\end{table}

\paragraph{Significance of the Tate result.}
The measured anomaly of $84 \pm 21$~ppm above the bare mass corresponds to
$92 \pm 21$~ppm above the BCS prediction.  The DFD value $\delta = 92.23$~ppm
lies within $0.01$~ppm of the central value of the Tate anomaly above BCS---well
within $0.4\sigma$.

\paragraph{Subsequent measurements.}
Hoang et al.\ (2019) performed high-precision measurements using a rotating
spherical-shell superconductor at ultra-low temperatures ($\sim 10^{-4}$~K),
extending the measurement to six superconducting materials~\cite{Hoang2019}.
For niobium, the gap to BCS predictions was reduced to 18~ppm but with a
dramatically improved accuracy of 2.1~ppm, making the disagreement \emph{more}
statistically significant ($\sim 8.6\sigma$) despite the smaller central value.

\paragraph{Tajmar experiments.}
Tajmar and de Matos (2006) proposed that the Tate anomaly could be explained by
a gravitomagnetic London moment~\cite{Tajmar2006}, and reported experimental
detection of anomalous acceleration fields near rotating superconductors.
However, subsequent high-precision measurements by Tajmar et al.\ (2022) using
a nano-Newton resolution double-pendulum thrust balance showed that earlier
anomalous force claims were likely electromagnetic artifacts~\cite{Tajmar2022}.

\subsection{Current Status}

The Cooper pair mass anomaly remains experimentally established and theoretically
unexplained within conventional physics.  The key facts are:
\begin{enumerate}[label=(\roman*)]
\item The anomaly is statistically significant ($> 4\sigma$ in Tate, $> 8\sigma$
in Hoang).
\item No conventional BCS extension has reproduced the measured value.
\item The anomaly persists across multiple superconducting materials.
\item DFD predicts $\delta = \sqrt{3}\,\alpha^2 = 92.23$~ppm from first
principles, matching the Tate measurement.
\end{enumerate}

%=============================================================================
\section{DFD Prediction: $\delta = \sqrt{3}\,\alpha^2$}\label{sec:dfd}
%=============================================================================

\subsection{Origin in the $A_5$ Microsector}

In DFD, the internal degrees of freedom of fermions are governed by the
\emph{microsector}---the channel Hilbert space
$\mathcal{H}_{\rm ch} = \mathbb{C}[A_5]$ of dimension $|A_5| = 60$.
The alternating group $A_5$ has five conjugacy classes:

\begin{center}
\renewcommand{\arraystretch}{1.3}
\begin{tabular}{lccc}
\toprule
Class & Representative & Size & Order \\
\midrule
1A & $e$ & 1 & 1 \\
2A & $(12)(34)$ & 15 & 2 \\
3A & $(123)$ & 20 & 3 \\
5A & $(12345)$ & 12 & 5 \\
5B & $(12354)$ & 12 & 5 \\
\bottomrule
\end{tabular}
\end{center}

The five irreducible representations of $A_5$ have dimensions $1, 3, 3', 4, 5$,
summing to $1^2 + 3^2 + 3^2 + 4^2 + 5^2 = 60 = |A_5|$ as required by the
Peter--Weyl decomposition.

\subsection{The Quintet Channel for Cooper Pairs}

A single unpaired electron occupies the fundamental representation $V^* \cong
\mathbb{C}^3$ of $A_5$ (one of the two 3-dimensional irreps, corresponding to
the three generations).  When two electrons form a Cooper pair, the composite
state lives in the symmetric tensor product:
\begin{equation}\label{eq:S2V}
S^2(V^*) = \mathbf{1} \oplus \mathbf{5}.
\end{equation}
The decomposition yields a singlet $|\mathbf{1}\rangle$ and a quintet
$|\mathbf{5}\rangle$.  The singlet channel reproduces the standard BCS pairing;
the quintet channel is \emph{new} and accessible only to paired states.

\subsection{Exchange Energy in the Quintet Channel}

The quintet $|\mathbf{5}\rangle$ is the 5-dimensional irreducible representation
of $A_5$.  Its exchange interaction with the macroscopic $\psi$-field contributes
an additional mass-energy to the Cooper pair.  The exchange amplitude is determined
by the structure constants of the $A_5$ Cayley graph propagator and the
$\psi$--matter coupling $\kappa = \alpha/4$:
\begin{equation}\label{eq:exchange}
\delta m_{\rm quintet} = 2m_e \times \sqrt{3}\,\alpha^2.
\end{equation}

The factor $\sqrt{3}$ arises from the $\mathbb{C}P^2$ geometry of the internal
space: the quintet channel connects through the diagonal of the $3 \times 3$
generation structure, producing a $\sqrt{3}$ geometric factor from the vertex
distance on $\mathbb{C}P^2$.  The $\alpha^2$ factor is the square of the
$\psi$--matter coupling, reflecting the second-order exchange process
(two vertices of the $\psi$--electron interaction).

\subsection{Numerical Evaluation}

\begin{align}
\alpha &= 7.2973525693 \times 10^{-3}, \\
\alpha^2 &= 5.3251 \times 10^{-5}, \\
\sqrt{3}\,\alpha^2 &= 9.2234 \times 10^{-5} = 92.23~\text{ppm}.
\end{align}

The predicted effective mass ratio is:
\begin{equation}
\frac{m^*}{2m_e}\bigg|_{\rm DFD} = 1 + \sqrt{3}\,\alpha^2 = 1.000\,092\,23.
\end{equation}

Comparing with the Tate measurement of $1.000\,084(21)$:
\begin{equation}
\frac{(m^*/2m_e)_{\rm DFD} - (m^*/2m_e)_{\rm Tate}}{\sigma_{\rm Tate}} =
\frac{0.000\,008\,23}{0.000\,021} = 0.39\sigma.
\end{equation}

This is a remarkably precise agreement from a \emph{zero-parameter} prediction.

%=============================================================================
\section{The Ross et al.\ (2024) E\"ot-Wash Constraint}\label{sec:ross}
%=============================================================================

\subsection{Experimental Overview}

Ross et al.\ (2024) performed the most precise test of the equivalence principle
for superconductors to date, using a cryogenic torsion balance at the University
of Washington~\cite{Ross2024}.  The key parameters of their experiment:

\begin{itemize}
\item \textbf{Test bodies:} Four 9.5~g niobium bodies and four 9.5~g copper
bodies (total pendulum mass 38~g), arranged on a torsion pendulum
\item \textbf{Suspension:} 30~$\mu$m tungsten fibre, 16.5~cm long
\item \textbf{Spring constant:} $\kappa = 8.7 \times 10^{-8}$~N$\cdot$m/rad
\item \textbf{Resonant frequency:} $\omega_0 = 2\pi \times 5.1$~mHz
\item \textbf{Temperature:} Inner shield at $\sim 5$~K; pendulum estimated
at $< 8$~K (well below $T_c = 9.25$~K for Nb)
\item \textbf{Data collection:} March 26 to April 22, 2024 (22 days usable)
\end{itemize}

\subsection{Results}

Two E\"otv\"os parameters were extracted:

\begin{tcolorbox}[title={Ross et al.\ 2024 Results}]
\begin{align}
\eta_{\rm Nb^*\text{-}Cu} &\leq 2.0 \times 10^{-9} \quad (95\%~\text{CL}),
\label{eq:ross_bulk} \\
\eta_{\rm CP\text{-}ee} &\leq 9.2 \times 10^{-4} \quad (95\%~\text{CL}).
\label{eq:ross_CP}
\end{align}
\end{tcolorbox}

The first bound (\ref{eq:ross_bulk}) constrains the \emph{bulk} equivalence
principle violation for superconducting niobium vs.\ copper.  The second bound
(\ref{eq:ross_CP}) constrains the E\"otv\"os parameter \emph{per Cooper pair}
vs.\ unpaired electrons.

\subsection{Cooper Pair Fraction}

Ross et al.\ used the two-fluid model to estimate the Cooper pair fraction:
\begin{equation}
\frac{N_{\rm CP}}{N_v} = 1 - \left(\frac{T}{T_c}\right)^4.
\end{equation}
With a conservative upper temperature estimate of $T = 8$~K and $T_c = 9$~K
for niobium:
\begin{equation}
f_{\rm CP} = \frac{N_{\rm CP}}{N_v} = 1 - \left(\frac{8}{9}\right)^4 = 0.38.
\end{equation}

\subsection{Constraint on $\lambda_{\rm quintet}$}

The DFD mass anomaly $\delta = 9.22 \times 10^{-5}$ predicts a per-Cooper-pair
E\"otv\"os parameter of:
\begin{equation}
\eta_{\rm CP}^{\rm DFD} = |1 - \lambda_{\rm quintet}| \times \delta.
\end{equation}

The Ross bound requires:
\begin{equation}
|1 - \lambda_{\rm quintet}| \times \delta \leq 9.2 \times 10^{-4},
\end{equation}
giving:
\begin{equation}\label{eq:lambda_bound}
\boxed{|1 - \lambda_{\rm quintet}| \leq \frac{9.2 \times 10^{-4}}{9.22 \times 10^{-5}} = 10.0}
\end{equation}

This means $\lambda_{\rm quintet} \in [-9, 11]$.  The Ross bound is
\emph{ten times too weak} to constrain $\lambda_{\rm quintet}$ to the region
near unity.  In particular, the propulsion-critical value
$\lambda_{\rm quintet} = 0$ is \textbf{completely allowed}.

\subsection{Why MICROSCOPE Cannot See This}

The MICROSCOPE satellite~\cite{MICROSCOPE2022} achieved $\eta \leq 1.5 \times 10^{-15}$
for titanium vs.\ platinum---but both are \emph{normal metals at room temperature}.
Neither material was superconducting during the measurement.  MICROSCOPE tested
only the singlet channel of unpaired electrons, which couples to $\psi$ normally.
The quintet channel is invisible to any experiment that does not involve a
superconducting condensate.

%=============================================================================
\section{Force Calculations}\label{sec:force}
%=============================================================================

\subsection{General Formula}

For a superconductor of total mass $M$ at temperature $T \ll T_c$:

\begin{itemize}
\item Cooper pair mass fraction: $f_{\rm CP} \approx 0.38$ (for Nb at $T < 8$~K)
\item Anomalous inertial mass: $\delta m_{\rm inert} = f_{\rm CP} \times M \times \delta$
\item Gravitational mass deficit (if $\lambda_{\rm quintet} \neq 1$):
\begin{equation}
\Delta m_{\rm grav} = f_{\rm CP} \times M \times \delta \times (1 - \lambda_{\rm quintet})
\end{equation}
\item Weight anomaly:
\begin{equation}\label{eq:weight}
\boxed{\Delta W = M \cdot g \cdot f_{\rm CP} \cdot \delta \cdot |1 - \lambda_{\rm quintet}|}
\end{equation}
\end{itemize}

\subsection{Numerical Results}

Table~\ref{tab:forces} presents the weight anomaly for various sample masses
and values of $\lambda_{\rm quintet}$, using $f_{\rm CP} = 0.38$,
$\delta = 9.223 \times 10^{-5}$, and $g = 9.807$~m/s$^2$.

\begin{table}[h]
\centering
\renewcommand{\arraystretch}{1.3}
\begin{tabular}{rcccc}
\toprule
\textbf{$M$ (kg)} & \multicolumn{4}{c}{\textbf{Weight anomaly $\Delta W$}} \\
\cmidrule{2-5}
& $\lambda = 0$ & $\lambda = 0.25$ & $\lambda = 0.5$ & $\lambda = 0.75$ \\
\midrule
0.1 & 0.034 mN & 0.026 mN & 0.017 mN & 0.009 mN \\
    & (3.5 mg$_f$) & (2.6 mg$_f$) & (1.8 mg$_f$) & (0.9 mg$_f$) \\[4pt]
1.0 & 0.344 mN & 0.258 mN & 0.172 mN & 0.086 mN \\
    & (35.0 mg$_f$) & (26.3 mg$_f$) & (17.5 mg$_f$) & (8.8 mg$_f$) \\[4pt]
5.0 & 1.719 mN & 1.289 mN & 0.859 mN & 0.430 mN \\
    & (175 mg$_f$) & (131 mg$_f$) & (87.6 mg$_f$) & (43.8 mg$_f$) \\[4pt]
\textbf{10.0} & \textbf{3.437 mN} & \textbf{2.578 mN} & \textbf{1.719 mN} & \textbf{0.859 mN} \\
    & \textbf{(350 mg$_f$)} & \textbf{(263 mg$_f$)} & \textbf{(175 mg$_f$)} & \textbf{(87.6 mg$_f$)} \\[4pt]
50.0 & 17.19 mN & 12.89 mN & 8.59 mN & 4.30 mN \\
    & (1.75 g$_f$) & (1.31 g$_f$) & (876 mg$_f$) & (438 mg$_f$) \\[4pt]
100.0 & 34.37 mN & 25.78 mN & 17.19 mN & 8.59 mN \\
    & (3.50 g$_f$) & (2.63 g$_f$) & (1.75 g$_f$) & (876 mg$_f$) \\
\bottomrule
\end{tabular}
\caption{Predicted weight anomaly for niobium samples of various masses and
gravitational coupling parameters.  Boldface: reference case of 10~kg.
The notation mg$_f$ denotes milligrams-force (weight equivalent).
For $\lambda_{\rm quintet} = 1$ (quintet mass gravitates normally), $\Delta W = 0$.}
\label{tab:forces}
\end{table}

\subsection{Reference Case: 10~kg Niobium}

For the reference case ($M = 10$~kg, $\lambda_{\rm quintet} = 0$):

\begin{tcolorbox}[title={Reference Calculation}]
\begin{align}
\text{Anomalous inertial mass} &= f_{\rm CP} \times M \times \delta \notag\\
&= 0.38 \times 10 \times 9.223 \times 10^{-5} \notag\\
&= 3.505 \times 10^{-4}~\text{kg} = 350.5~\text{mg} \\[6pt]
\text{Weight anomaly} &= M \cdot g \cdot f_{\rm CP} \cdot \delta \notag\\
&= 10 \times 9.807 \times 0.38 \times 9.223 \times 10^{-5} \notag\\
&= 3.437 \times 10^{-3}~\text{N} = 3.437~\text{mN}
\end{align}
\end{tcolorbox}

This corresponds to a weight change of 350~mg---easily detectable on a standard
analytical balance with 0.01~mg resolution (signal-to-noise ratio $35{,}000\!:\!1$).

Even for $\lambda_{\rm quintet} = 0.99$ (1\% violation), the weight anomaly
for 10~kg Nb is 3.5~mg, still detectable at $350\!:\!1$ SNR.

\subsection{Physical Properties of Niobium Samples}

\begin{itemize}
\item Niobium density: $\rho_{\rm Nb} = 8{,}570$~kg/m$^3$
\item Critical temperature: $T_c = 9.25$~K (Type~II superconductor)
\item Volume of 10~kg sample: 1{,}167~cm$^3$ (cube of side 10.5~cm)
\item Niobium is commercially available in large quantities
\item Cost of 10~kg Nb: approximately \$3{,}000--\$5{,}000
\end{itemize}

%=============================================================================
\section{Experimental Design}\label{sec:experiment}
%=============================================================================

\subsection{Concept}

The experiment is conceptually simple: weigh a niobium sample \emph{above} and
\emph{below} $T_c$ and look for a weight change that correlates with the
superconducting transition.

\begin{tcolorbox}[title={Experimental Protocol}]
\begin{enumerate}
\item Cool Nb sample to $T_0 \approx 4$~K (well below $T_c = 9.25$~K)
\item Measure weight $W_{\rm SC}$ (superconducting state)
\item Heat to $T_1 \approx 12$~K (above $T_c$)
\item Measure weight $W_N$ (normal state)
\item Repeat Steps 1--4 for $N \geq 100$ thermal cycles
\item Search for systematic weight difference $\Delta W = W_N - W_{\rm SC}$
correlated with the superconducting transition
\end{enumerate}
\end{tcolorbox}

\paragraph{Expected signal.}
For 10~kg Nb with $\lambda_{\rm quintet} = 0$:
\begin{equation}
\Delta W = W_N - W_{\rm SC} = +3.437~\text{mN} = +350~\text{mg}_f.
\end{equation}
The superconducting sample is \emph{lighter} (some of its mass doesn't gravitate),
so $W_N > W_{\rm SC}$.

\subsection{Balance Requirements}

\begin{table}[h]
\centering
\renewcommand{\arraystretch}{1.3}
\begin{tabular}{lcc}
\toprule
\textbf{Balance Type} & \textbf{Resolution} & \textbf{SNR (10 kg, $\lambda=0$)} \\
\midrule
Industrial (0.1~g) & 100 mg & 3.5:1 \\
Precision (0.01~g) & 10 mg & 35:1 \\
Analytical (0.01~mg) & 0.01 mg & 35,000:1 \\
Microbalance (0.001~mg) & 0.001 mg & 350,000:1 \\
\bottomrule
\end{tabular}
\caption{Signal-to-noise ratio for different balance types.  Even a basic
precision balance provides adequate SNR.  An analytical balance provides
overwhelming sensitivity.}
\label{tab:balance}
\end{table}

In practice, a cryogenic mass comparator must be purpose-built (no commercial
analytical balance operates at 4~K), but the required force sensitivity of
$\sim 10~\mu$N is well within the capability of:
\begin{itemize}
\item Capacitance-bridge force sensors
\item Torsion balances (as used by Ross et al.)
\item Piezoelectric load cells
\item SQUID-based force measurement
\end{itemize}

\subsection{Systematic Error Budget}\label{sec:systematics}

\subsubsection{Meissner Effect (Magnetic Forces)}

When niobium transitions to the superconducting state in Earth's magnetic field
($B_\oplus \approx 50~\mu$T), it expels magnetic flux (Meissner effect).  This
produces a force in any field gradient.

\paragraph{Estimate.}
The magnetic force on a perfect diamagnet of volume $V$ in a field gradient is:
\begin{equation}
F_{\rm mag} = \frac{|\chi| V B \, (\nabla B)}{\mu_0}.
\end{equation}
With $V = 1.17 \times 10^{-3}$~m$^3$, $B = 5 \times 10^{-5}$~T, and
Earth's field gradient $\nabla B \approx 3 \times 10^{-6}$~T/m:
\begin{equation}
F_{\rm mag} \approx 1.4 \times 10^{-7}~\text{N} = 0.14~\mu\text{N}.
\end{equation}

This is \textbf{24,000 times smaller} than the predicted weight anomaly
(3.44~mN).  Nevertheless, it should be controlled by:
\begin{itemize}
\item Magnetic shielding (mu-metal + superconducting Pb shield) to reduce
$B$ at the sample to $< 1~\mu$T
\item Placing the sample at the centre of a Helmholtz coil null
\item Using a geometry where the field gradient averages to zero over the sample
\end{itemize}

\subsubsection{Buoyancy Changes}

If the experiment uses helium exchange gas for thermal contact, the gas density
changes between 4~K and 12~K, producing a buoyancy shift.

\paragraph{Estimate.}
Helium density at 4~K (1~atm): $\rho_{\rm He} \approx 0.16$~kg/m$^3$.
Helium density at 10~K (1~atm): $\rho_{\rm He} \approx 0.13$~kg/m$^3$.
\begin{equation}
\Delta F_{\rm buoy} = (\rho_{4K} - \rho_{10K}) \cdot V \cdot g
= 0.03 \times 1.17 \times 10^{-3} \times 9.81 = 3.4 \times 10^{-4}~\text{N}
= 0.34~\text{mN}.
\end{equation}

This is \textbf{10\% of the predicted signal}---a significant systematic.

\paragraph{Mitigation.}
\begin{itemize}
\item \textbf{Operate in vacuum} ($< 10^{-3}$~mbar): eliminates buoyancy entirely.
Thermal contact is maintained by mechanical conduction through the suspension.
\item If exchange gas is needed: use precise $T$--$P$ monitoring and compute the
buoyancy correction; or use a dummy sample (copper, same volume) as a
differential reference.
\end{itemize}

\subsubsection{Thermal Contraction}

Niobium's linear thermal expansion coefficient is $\sim 7 \times 10^{-6}$~K$^{-1}$
at room temperature, dropping to $\sim 10^{-7}$~K$^{-1}$ below 20~K.

\paragraph{Estimate.}
Volume change from 12~K to 4~K:
\begin{equation}
\frac{\Delta V}{V} \approx 3 \times 10^{-7} \times 8 = 2.4 \times 10^{-6}.
\end{equation}
This affects buoyancy by $\sim 10^{-6}$~mN---completely negligible.

\subsubsection{Electromagnetic Artifacts}

The superconducting transition involves redistribution of screening currents,
which can produce forces through interaction with stray fields.

\paragraph{Mitigation.}
\begin{itemize}
\item Use DC SQUID monitoring to track any persistent currents
\item Cycle through $T_c$ both up and down; artifacts from trapped flux
produce different signatures on warming vs.\ cooling
\item Use multiple samples with different geometries (same mass, different
shape) to separate geometric-dependent EM artifacts from the mass-dependent
gravitational signal
\end{itemize}

\subsubsection{Summary of Systematics}

\begin{table}[h]
\centering
\renewcommand{\arraystretch}{1.3}
\begin{tabular}{lcc}
\toprule
\textbf{Systematic} & \textbf{Magnitude} & \textbf{Ratio to Signal} \\
\midrule
Meissner force (unshielded) & $0.14~\mu$N & $4 \times 10^{-5}$ \\
Meissner force (shielded to 1~$\mu$T) & $0.003~\mu$N & $8 \times 10^{-7}$ \\
Buoyancy (He gas, 4$\to$12 K) & 0.34 mN & 0.10 \\
Buoyancy (vacuum) & $< 10^{-6}$ mN & $< 3 \times 10^{-7}$ \\
Thermal contraction & $10^{-6}$ mN & $3 \times 10^{-7}$ \\
EM artifacts & variable & controlled by protocol \\
\midrule
\textbf{Predicted signal ($\lambda = 0$)} & \textbf{3.44 mN} & \textbf{1.0} \\
\bottomrule
\end{tabular}
\caption{Systematic error budget for the 10~kg niobium experiment.  All
systematics are controllable to levels far below the predicted signal.}
\label{tab:systematics}
\end{table}

\subsection{Thermal Cycling Protocol}

The key experimental strategy is \emph{repeated thermal cycling}:

\begin{enumerate}
\item Stabilize at $T_1 = 12$~K for 30~minutes; record weight.
\item Cool through $T_c$ to $T_0 = 4$~K; stabilize for 30~minutes; record weight.
\item Heat through $T_c$ to $T_1 = 12$~K; stabilize for 30~minutes; record weight.
\item Repeat for $N \geq 100$ complete cycles.
\end{enumerate}

The expected signal is a \emph{step function} in weight at $T_c = 9.25$~K,
switching on and off with each cycle.  This is readily distinguished from
slow drifts, linear trends, or random noise by:
\begin{itemize}
\item Lock-in analysis at the cycling frequency
\item Phase correlation with independently measured $T_c$ transition
(monitored by resistivity or AC susceptibility)
\item Statistical analysis of the weight-change distribution
\end{itemize}

With 100 cycles and even a modest single-cycle SNR of 10:1, the accumulated
significance exceeds $100\sigma$.

\subsection{Cost and Feasibility}

\begin{table}[h]
\centering
\renewcommand{\arraystretch}{1.3}
\begin{tabular}{lr}
\toprule
\textbf{Item} & \textbf{Estimated Cost} \\
\midrule
10 kg niobium sample & \$5,000 \\
Cryostat (closed-cycle, 4 K) & \$25,000 \\
Force sensor / load cell & \$5,000 \\
Magnetic shielding & \$3,000 \\
Temperature control \& monitoring & \$5,000 \\
Data acquisition & \$3,000 \\
\midrule
\textbf{Total} & \textbf{\$46,000} \\
\bottomrule
\end{tabular}
\caption{Estimated cost for the Cooper pair mass anomaly weight test.
Many university cryogenics labs already possess the cryostat and DAQ,
reducing the marginal cost to $\sim$\$15,000.}
\label{tab:cost}
\end{table}

This experiment can be performed in \textbf{any university cryogenics laboratory}
with a 4~K cryostat.  No particle accelerator, no space mission, no high-$Q$
cavity---just a chunk of niobium and a good scale.

%=============================================================================
\section{Why This Mechanism Is Fundamentally Different}\label{sec:comparison}
%=============================================================================

\subsection{The Problem with Previous Propulsion Proposals}

All previous attempts at ``gravitational propulsion'' from electromagnetic
energy encounter a fundamental barrier.  The coupling between electromagnetic
energy and spacetime curvature (or in DFD, the $\psi$-field) goes through the
gravitational coupling $G/c^4$:
\begin{equation}
\frac{2G}{c^4} = 1.65 \times 10^{-44}~\text{m/J}.
\end{equation}

To produce 1~N of force by \emph{creating} a $\psi$-gradient from EM energy
requires storing $\sim 10^{44}$~J---the rest-mass energy of Jupiter.

Previous proposals include:
\begin{itemize}
\item \textbf{Psi-bubble propulsion:} Creates $\psi$-gradient using high-$Q$ SRF
cavities and rotating EM fields.  Achievable forces: $\sim 10^{-35}$~N for
laboratory-scale systems.
\item \textbf{Podkletnov rotating superconductor:} Claimed weight reduction above
a rotating YBCO disk.  Never independently replicated; likely electromagnetic
artifacts.
\item \textbf{Tajmar gravitomagnetic London moment:} Reported anomalous
accelerations near rotating superconductors.  Later shown by Tajmar's own
group to be electromagnetic artifacts at nano-Newton precision.
\end{itemize}

\subsection{The Cooper Pair Mechanism Is Different}

The Cooper pair mass anomaly mechanism does \textbf{not} attempt to create a
$\psi$-gradient.  Instead, it exploits a pre-existing feature of the DFD
matter--$\psi$ coupling:

\begin{tcolorbox}[title={Fundamental Distinction}]
\textbf{Previous approaches:} Supply EM energy $\to$ create $\psi$-gradient
$\to$ produce force.\\
$\Rightarrow$ Blocked by the $2G/c^4$ wall.

\textbf{Cooper pair mechanism:} Modify the \emph{coupling} between existing
matter and the \emph{existing} $\psi$-field (Earth's gravitational field) by
switching the quantum state of the condensate.\\
$\Rightarrow$ No energy barrier.  The anomalous mass $\delta = 92$~ppm is
\emph{already there} (Tate measured it in 1989).
\end{tcolorbox}

The physics is analogous to the difference between:
\begin{itemize}
\item Creating an electric field from scratch (requires charge separation, costs energy)
\item Changing the dielectric constant of a material in an existing field (phase
transition, comparatively trivial)
\end{itemize}

The superconducting transition is the ``dielectric switch'' for gravity.  It
costs only the condensation energy ($\sim$~meV per Cooper pair), but it modifies
how $\sim 92$~ppm of the condensate's mass-energy couples to $\psi$.

\subsection{Comparison Table}

\begin{table}[h]
\centering
\renewcommand{\arraystretch}{1.3}
\begin{tabular}{lccc}
\toprule
\textbf{Mechanism} & \textbf{Force (10 kg)} & \textbf{Cost} & \textbf{Complexity} \\
\midrule
$\psi$-bubble (SRF cavity) & $\sim 10^{-35}$ N & \$10M+ & Extreme \\
Podkletnov (rotating SC) & unconfirmed & \$1M+ & High \\
Tajmar (gravitomagnetic) & null result & \$500K+ & High \\
\textbf{Cooper pair anomaly} & \textbf{3.44 mN} & \textbf{\$50K} & \textbf{Low} \\
\bottomrule
\end{tabular}
\caption{Comparison of propulsion mechanisms.  The Cooper pair anomaly produces
a force $\sim 10^{32}$ times larger than the $\psi$-bubble mechanism, at
1/200th the cost, with a clear experimental signature.  The only caveat is
the unknown value of $\lambda_{\rm quintet}$.}
\label{tab:comparison}
\end{table}

%=============================================================================
\section{Theoretical Question: What Determines $\lambda_{\rm quintet}$?}
\label{sec:theory}
%=============================================================================

\subsection{The Central Question}

The entire propulsion viability hinges on one theoretical question:

\begin{tcolorbox}[colframe=red!75!black, colback=red!5!white,
title={The $\lambda_{\rm quintet}$ Question}]
Does the quintet exchange energy in $S^2(V^*)$ contribute to the
stress-energy tensor $T_{\mu\nu}$ that sources the $\psi$-field?

\begin{itemize}
\item $\lambda_{\rm quintet} = 1$: Yes $\to$ EP preserved, no propulsion
\item $\lambda_{\rm quintet} = 0$: No $\to$ EP violated for Cooper pairs, propulsion viable
\item $0 < \lambda_{\rm quintet} < 1$: Partial $\to$ reduced but nonzero force
\end{itemize}
\end{tcolorbox}

\subsection{DFD Action Principle Considerations}

In the DFD framework, matter couples to $\psi$ through the frame stiffness
action:
\begin{equation}
\mathcal{L}_{\rm matter} = -\kappa_0\,\psi \cdot S[\rho],
\end{equation}
where $\rho$ is the internal density matrix and $S[\rho]$ is the von Neumann
entropy of the microsector.  The partition function $Z(S^3)$ encodes the sum
over all microsector states.

The quintet channel in $S^2(V^*)$ arises from the \emph{pairing} of two
electrons within the $V^*$ representation.  Its contribution to $S[\rho]$
depends on whether the bridge map
\begin{equation}
\mathcal{B}: \mathcal{H}_{\rm ch} \to \mathcal{H}_{\rm macro}
\end{equation}
preserves or projects out the quintet component when mapping microsector
states to macroscopic observables that source $\psi$.

\subsection{Three Scenarios}

\paragraph{Scenario A: Universal coupling ($\lambda_{\rm quintet} = 1$).}
All internal energy gravitates.  The bridge map is the identity on the
energy eigenspace.  The EP is exactly preserved for Cooper pairs.  The Tate
anomaly exists in the inertial mass but is exactly matched in the gravitational
mass.  No propulsion.

\paragraph{Scenario B: Projected coupling ($\lambda_{\rm quintet} = 0$).}
The bridge map projects onto the singlet channel $|\mathbf{1}\rangle \in S^2(V^*)$,
discarding the quintet $|\mathbf{5}\rangle$.  Physically, this means the
quintet exchange energy is ``internal'' to the microsector and does not appear
in the macroscopic stress-energy.  The EP is violated for Cooper pairs by
exactly $\delta = 92$~ppm.  Full propulsion force.

This scenario has a natural interpretation: the quintet channel represents an
\emph{internal} degree of freedom of the $A_5$ microsector that does not couple
to the macroscopic $\psi$-field, analogous to how the colour degree of freedom
of quarks does not produce macroscopic colour fields due to confinement.

\paragraph{Scenario C: Partial coupling ($0 < \lambda_{\rm quintet} < 1$).}
The bridge map partially transmits the quintet energy.  This could arise from
finite-temperature or finite-coherence-length corrections to the bridge map.
Produces a reduced but nonzero force proportional to $(1 - \lambda_{\rm quintet})$.

\subsection{What Theory Predicts (Open Question)}

Determining $\lambda_{\rm quintet}$ from first principles requires computing the
bridge map $\mathcal{B}$ for the specific case of the quintet representation.
This involves:

\begin{enumerate}
\item The structure of $S^2(V^*)$ as a representation of $A_5$
\item The action of the Cayley graph propagator on the quintet subspace
\item Whether the $\psi$-field sources from the full microsector partition
function or only from its projection onto specific representation channels
\item The role of gauge invariance: if the quintet transforms nontrivially
under the residual gauge group, its expectation value may vanish in the
macroscopic limit (analogous to colour confinement)
\end{enumerate}

This calculation is the \textbf{most important open theoretical problem in the
DFD propulsion programme}.  However, it is also a question that can be settled
\emph{experimentally} far more quickly and cheaply than theoretically, using the
weight-change test described in Section~\ref{sec:experiment}.

%=============================================================================
\section{Implications and Next Steps}\label{sec:implications}
%=============================================================================

\subsection{If $\lambda_{\rm quintet} \neq 1$ (Propulsion Works)}

The implications are transformative:

\begin{enumerate}
\item \textbf{Immediate:} The equivalence principle is violated for
superconducting condensates.  This is the first macroscopic EP violation and
constitutes a discovery of the highest magnitude in fundamental physics.

\item \textbf{Engineering:} A passive gravitational shielding effect of
$\delta \cdot f_{\rm CP} = 35$~ppm is available by cooling niobium below $T_c$.
No energy input beyond cryogenic cooling is needed.

\item \textbf{Propulsion:} For a 1000~kg Nb payload at 4~K:
\begin{equation}
F = 1000 \times 9.81 \times 0.38 \times 9.22 \times 10^{-5} = 0.344~\text{N}.
\end{equation}
This is not enough for launch but is significant for station-keeping,
attitude control, and long-duration deep-space missions.

\item \textbf{Scaling:} Higher-$T_c$ superconductors with larger condensate
fractions could enhance the effect.  If the anomaly scales with $T_c$ (via
increased pairing interaction), materials like MgB$_2$ ($T_c = 39$~K) or
cuprates ($T_c \sim 90$~K) may show larger effects.

\item \textbf{Amplification:} If the anomalous mass can be modulated
(e.g., by driving the condensate through $T_c$ at high frequency), the
resulting oscillating gravitational mass could be used for gravitational
wave generation or resonant coupling to external masses.
\end{enumerate}

\subsection{If $\lambda_{\rm quintet} = 1$ (No Propulsion)}

Even a null result is highly valuable:

\begin{enumerate}
\item It confirms the EP for superconducting condensates at the 0.01~mg level---a
new precision frontier.
\item Combined with the DFD prediction of $\delta$, it constrains the bridge
map to be representation-preserving, which is important theoretical information.
\item It motivates the search for other channels where $\lambda \neq 1$ might hold.
\end{enumerate}

\subsection{Experimental Roadmap}

\begin{enumerate}
\item \textbf{Phase 1 (6 months, \$50K):} Niobium weight-change test at a
university cryogenics lab.  Goal: detect or bound $\Delta W$ at the 0.1~mg level.

\item \textbf{Phase 2 (12 months, \$200K):} If signal detected, replicate with:
\begin{itemize}
\item Multiple Nb samples of different masses (1, 5, 10, 50~kg)
\item Verify $\Delta W \propto M$ scaling
\item Test temperature dependence: $\Delta W(T)$ should follow the BCS gap
function $\Delta(T)$ below $T_c$
\end{itemize}

\item \textbf{Phase 3 (18 months, \$500K):} Extend to other superconductors
(Pb, Sn, Al, MgB$_2$, YBCO) to test material dependence.

\item \textbf{Phase 4 (24 months):} If the effect is confirmed across materials,
begin engineering optimization for maximum force-to-mass ratio.
\end{enumerate}

%=============================================================================
\section{Conclusions}\label{sec:conclusions}
%=============================================================================

\begin{enumerate}
\item The Tate Cooper pair mass anomaly ($92 \pm 21$~ppm above BCS) has been
unexplained for 36 years.  DFD predicts \emph{exactly} $\delta = \sqrt{3}\,\alpha^2
= 92.23$~ppm from the quintet channel of the $A_5$ microsector.

\item The Ross et al.\ (2024) E\"ot-Wash bound $\eta_{\rm CP} \leq 9.2 \times 10^{-4}$
is ten times too weak to constrain the gravitational coupling parameter
$\lambda_{\rm quintet}$.  The propulsion-relevant value $\lambda_{\rm quintet} = 0$
is completely allowed.

\item For a 10~kg niobium sample with $\lambda_{\rm quintet} = 0$, the predicted
weight anomaly is 3.44~mN (350~mg$_f$), detectable at $35{,}000\!:\!1$ SNR on
a 0.01~mg balance.

\item All systematic errors (Meissner effect, buoyancy, thermal contraction,
EM artifacts) are controllable to levels far below the predicted signal.

\item The experiment can be performed in any university cryogenics laboratory
for less than \$50{,}000.  This is the cheapest and simplest test of any
proposed propulsion mechanism by many orders of magnitude.

\item Unlike all previous propulsion proposals, this mechanism does not create
$\psi$-gradients from EM energy (and thus is not blocked by the $2G/c^4$ wall).
It modifies the \emph{coupling} between existing matter and the existing
gravitational field via a quantum phase transition.

\item The theoretical question of $\lambda_{\rm quintet}$ can be settled
experimentally far more quickly than theoretically.  \textbf{This experiment
should be the highest priority in the DFD programme.}
\end{enumerate}

%=============================================================================
\begin{thebibliography}{99}

\bibitem{Tate1989}
J.~Tate, B.~Cabrera, S.B.~Felch, and J.T.~Anderson,
``Precise determination of the Cooper-pair mass,''
\textit{Phys.\ Rev.\ Lett.}\ \textbf{62}, 845--848 (1989).

\bibitem{Tate1990}
J.~Tate, B.~Cabrera, S.B.~Felch, and J.T.~Anderson,
``Determination of the Cooper-pair mass in niobium,''
\textit{Phys.\ Rev.\ B}\ \textbf{42}, 7885--7893 (1990).

\bibitem{Hoang2019}
C.~Hoang, M.~Bajraktari, and M.~Tajmar,
``High-precision measurement of Cooper-pair mass using rotating
spherical-shell superconductor,''
\textit{Mater.\ Lett.}\ \textbf{262}, 127178 (2020).

\bibitem{Ross2024}
M.P.~Ross, S.M.~Fleischer, I.A.~Paulson, P.~Lamb, B.M.~Iritani,
E.G.~Adelberger, C.A.~Hagedorn, K.~Venkateswara, C.~Gettings,
E.A.~Shaw, S.K.~Apple, and J.H.~Gundlach,
``Test of the equivalence principle for superconductors,''
arXiv:2407.21232 (2024).

\bibitem{Tajmar2006}
M.~Tajmar and C.J.~de Matos,
``Gravitomagnetic fields in rotating superconductors to solve
Tate's Cooper pair mass anomaly,''
\textit{AIP Conf.\ Proc.}\ \textbf{813}, 1415 (2006);
arXiv:gr-qc/0607086.

\bibitem{Tajmar2022}
M.~Tajmar, O.~Neunzig, and M.~K\"o{\ss}ling,
``Measurement of anomalous forces from a Cooper-pair current in
high-$T_c$ superconductors with nano-Newton precision,''
\textit{Front.\ Phys.}\ \textbf{10}, 892215 (2022).

\bibitem{DeWitt1966}
B.S.~DeWitt,
``Superconductors and gravitational drag,''
\textit{Phys.\ Rev.\ Lett.}\ \textbf{16}, 1092 (1966).

\bibitem{MICROSCOPE2022}
P.~Touboul \textit{et al.}\ (MICROSCOPE Collaboration),
``MICROSCOPE mission: Final results of the test of the equivalence
principle,''
\textit{Phys.\ Rev.\ Lett.}\ \textbf{129}, 121102 (2022).

\bibitem{DFD2026}
G.T.~Alcock,
\textit{Density Field Dynamics: A Complete Unified Theory}, v3.2 (2026).

\end{thebibliography}

%=============================================================================
\appendix
\section{Detailed Numerical Calculations}\label{app:numerics}

\subsection{Physical Constants Used}

\begin{align}
\alpha &= 1/137.035\,999\,084 = 7.2974 \times 10^{-3} \\
m_e &= 9.1094 \times 10^{-31}~\text{kg} \\
g &= 9.8067~\text{m/s}^2 \\
\rho_{\rm Nb} &= 8{,}570~\text{kg/m}^3 \\
T_{c,\rm Nb} &= 9.25~\text{K}
\end{align}

\subsection{DFD Anomaly Derivation Summary}

\begin{align}
\alpha^2 &= 5.3251 \times 10^{-5} \\
\sqrt{3} &= 1.7321 \\
\delta = \sqrt{3}\,\alpha^2 &= 9.2234 \times 10^{-5} = 92.23~\text{ppm}
\end{align}

\subsection{Ross Bound Analysis}

\begin{align}
\eta_{\rm CP} &\leq 9.2 \times 10^{-4} \\
\delta_{\rm DFD} &= 9.22 \times 10^{-5} \\
\frac{\eta_{\rm CP}}{\delta_{\rm DFD}} &= \frac{9.2 \times 10^{-4}}{9.22 \times 10^{-5}} = 9.98 \approx 10.0 \\
|1 - \lambda_{\rm quintet}| &\leq 10.0
\end{align}

\subsection{Force Calculation Verification}

For $M = 10$~kg, $\lambda_{\rm quintet} = 0$:
\begin{align}
\Delta W &= M \cdot g \cdot f_{\rm CP} \cdot \delta \cdot |1 - \lambda| \\
&= 10 \times 9.807 \times 0.38 \times 9.223 \times 10^{-5} \times 1.0 \\
&= 10 \times 9.807 \times 3.505 \times 10^{-5} \\
&= 10 \times 3.437 \times 10^{-4} \\
&= 3.437 \times 10^{-3}~\text{N} = 3.437~\text{mN}
\end{align}

Weight equivalent:
\begin{equation}
\Delta m_f = \frac{\Delta W}{g} = \frac{3.437 \times 10^{-3}}{9.807}
= 3.505 \times 10^{-4}~\text{kg} = 350.5~\text{mg}
\end{equation}

\end{document}
